\documentclass[12pt]{article}
\usepackage[utf8]{inputenc}
\usepackage[T1]{fontenc}
\usepackage{hyperref}
\title{Tuyuca Language Resources for K-12 Curriculum Development}
\author{Compiled by Claude Sonnet 5 (claude-sonnet-5) for Andr\'as Kornai}
\date{}
\begin{document}
\maketitle

\section{Introduction}
Tuyuca (ISO 639-3 tue; Glottocode tuyu1244) is an Eastern Tucanoan language spoken along the Papuri, Inambú and Tiquié rivers in Colombia and Brazil. Speaker estimates in the searches were around 1,400. It is known for its obligatory evidentiality system. The searches returned mostly linguistic materials. No school curriculum materials were found.

\section{Primary Educational Resources}
A learner-oriented Spanish--Tuyuca and Tuyuca--Spanish dictionary, \textit{Mi primer diccionario}, was published by the Instituto Lingüístico de Verano. Catalogue entries disagree on whether Barnes or C. Tamayo L. is the author. No URL was seen for it. No primers or textbooks were found.

\section{Linguistic and Reference Materials}
\begin{itemize}
\item Wikipedia overview of Tuyuca: \url{https://en.wikipedia.org/wiki/Tuyuca_language}
\item Janet Barnes 1992, \textit{Diccionario tuyuca-español} (3,749 entries, Oxford Text Archive), as catalogued by Glottolog: \url{https://glottolog.org/resource/reference/id/75836}
\item Janet Barnes 2000, Tuyuca dictionary in the Intercontinental Dictionary Series, as catalogued by Glottolog: \url{https://glottolog.org/resource/reference/id/20718}
\item ASJP word-list entry: \url{https://asjp.clld.org/languages/TUYUCA_2}
\item Channumerals page (numerals, not checked in detail): \url{https://lingweb.eva.mpg.de/channumerals/Tuyuca.htm}
\end{itemize}
Barnes's work on evidentiality and Barnes and Silzer's phonology (1976) were mentioned in the search summary, but no direct URLs for them were seen.

\section{Religious and Community Resources}
None found.

\section{Digital and Audio Resources}
No audio, video or software resources were found.

\section{Curriculum Development Context}
Tuyuca is used by every generation, so a mother-tongue curriculum is plausible. The searches found no school programme and no teaching materials beyond the dictionaries. A curriculum would need new materials developed with the community, and the sources above would serve as a lexical and phonological base.
\end{document}